
%%%%%%%%%%%%%%%%%%%%%%%%%%%%%%%%%%%%%%%%%%%%%%%%%%%%%%%%%%%%%%%%%%%%%%%
%% ARCHIVO: flowcharts/excess_enthalpy_submodule.tex
%% Submódulo de entalpía en exceso
%%%%%%%%%%%%%%%%%%%%%%%%%%%%%%%%%%%%%%%%%%%%%%%%%%%%%%%%%%%%%%%%%%%%%%%

\begin{figure}[H]
\centering
\begin{tikzpicture}[
    node distance=0.6cm,
    font=\small\sffamily,
    every node/.style={align=center}
]

% === INICIO SUBMÓDULO ===
\node (start_sub) [startstop] {INICIO\\Submódulo $\Delta H_{\text{exceso}}$};

% === ENTRADAS ===
\node (inputs_sub) [doc, below=0.5cm of start_sub, text width=4.5cm] {
    \textbf{Entradas:}\\[1pt]
    $T_{bs}$, $P_{\text{tot}}$, $P_v$\\
    $Y$, $z$, $r$\\
    Propiedades: $T_c$, $P_c$, $\omega$
};

% === DECISIÓN: TIPO DE MODELO ===
\node (model_choice) [decision, below=0.7cm of inputs_sub] {
    ¿Usar modelo\\simplificado?
};

% === OPCIÓN 1: MODELO SIMPLIFICADO (IDEAL) ===
\node (simple_model) [process, right=2cm of model_choice, text width=4cm] {
    \textbf{Modelo ideal:}\\[1pt]
    $\Delta H_{\text{exceso}} = 0$\\[1pt]
    {\scriptsize Válido para gases ideales}
};

% === OPCIÓN 2: MODELO COMPLETO ===
\node (complex_model) [process, left=4cm of model_choice, text width=4.5cm] {
    Calcular propiedades reducidas:\\[1pt]
    $\displaystyle T_r = \frac{T_{abs}}{T_c}$\\[1pt]
    $\displaystyle P_r = \frac{P_{\text{tot}}}{P_c}$
};

% === CÁLCULO COEFICIENTE VIRIAL ===
\node (calc_virial) [process, below=0.7cm of complex_model, text width=5cm] {
    Calcular coeficiente virial $B$:\\[1pt]
    $\displaystyle B = \frac{RT_c}{P_c}(B^{(0)} + \omega B^{(1)})$\\[1pt]
    {\scriptsize Usar correlaciones de Pitzer}
};

% === CÁLCULO dB/dT ===
\node (calc_dbdt) [process, below=0.6cm of calc_virial, text width=5.5cm] {
    Calcular derivada del coeficiente virial:\\[1pt]
    $\displaystyle \frac{dB}{dT} = \frac{R}{P_c}\left(\frac{dB^{(0)}}{dT_r} + \omega\frac{dB^{(1)}}{dT_r}\right)$
};

% === CORRELACIONES PITZER ===
\node (pitzer) [process, below=0.6cm of calc_dbdt, text width=6cm] {
    \textbf{Correlaciones de Pitzer:}\\[1pt]
    $\displaystyle B^{(0)} = 0.083 - \frac{0.422}{T_r^{1.6}}$\\[1pt]
    $\displaystyle B^{(1)} = 0.139 - \frac{0.172}{T_r^{4.2}}$\\[1pt]
    $\displaystyle \frac{dB^{(0)}}{dT_r} = \frac{0.675}{T_r^{2.6}}$, \quad
    $\displaystyle \frac{dB^{(1)}}{dT_r} = \frac{0.722}{T_r^{5.2}}$
};

% === CÁLCULO ENTALPÍA RESIDUAL GAS SECO ===
\node (calc_hr_gs) [process, below=0.6cm of pitzer, text width=6cm] {
    Entalpía residual del gas seco:\\[1pt]
    $\displaystyle H_{\text{gs}}^{\text{res}} = P_{\text{tot}}\left(B - T\frac{dB}{dT}\right)$
};

% === CÁLCULO ENTALPÍA RESIDUAL VAPOR ===
\node (calc_hr_v) [process, below=0.6cm of calc_hr_gs, text width=6cm] {
    Entalpía residual del vapor:\\[1pt]
    $\displaystyle H_{v}^{\text{res}} = P_v\left(B_v - T\frac{dB_v}{dT}\right)$\\[1pt]
    {\scriptsize Usar $B_v$ con propiedades del vapor}
};

% === CÁLCULO ENTALPÍA MEZCLA ===
\node (calc_mixing) [process, below=0.6cm of calc_hr_v, text width=6.5cm] {
    Entalpía de mezcla no ideal:\\[1pt]
    $\displaystyle \Delta H_{\text{mezcla}} = P_{\text{tot}} y_v y_{gs} \delta_{vgs}$\\[1pt]
    {\scriptsize $\delta_{vgs}$ = parámetro de interacción binaria}\\
    {\scriptsize $y_v = P_v/P_{\text{tot}}$, \quad $y_{gs} = 1 - y_v$}
};

% === CÁLCULO FINAL ===
\node (calc_excess) [process, below=0.6cm of calc_mixing, text width=6.5cm] {
    \textbf{Entalpía en exceso total:}\\[2pt]
    $\displaystyle \Delta H_{\text{exceso}} = H_{\text{gs}}^{\text{res}} + Y \cdot H_{v}^{\text{res}} + \Delta H_{\text{mezcla}}$
};

% === PUNTO DE CONVERGENCIA ===
\node (converge) [process, below=3cm of simple_model, text width=4cm] {
    Valor de $\Delta H_{\text{exceso}}$\\
    determinado
};

% === SALIDA ===
\node (output_sub) [doc, right=0.6cm of converge, text width=4cm] {
    \textbf{Salida:}\\[1pt]
    $\Delta H_{\text{exceso}}$
};

\node (end_sub) [startstop, right=0.5cm of output_sub] {FIN\\Submódulo};

% === CONEXIONES ===
\draw [arrow] (start_sub) -- (inputs_sub);
\draw [arrow] (inputs_sub) -- (model_choice);

% Ramificación del modelo
\draw [arrow_yes] (model_choice) -- node[label_node, above] {\isTrue} (simple_model);
\draw [arrow_no] (model_choice) -- node[label_node, above] {\isFalse} (complex_model);

% Flujo modelo complejo
\draw [arrow] (complex_model) -- (calc_virial);
\draw [arrow] (calc_virial) -- (calc_dbdt);
\draw [arrow] (calc_dbdt) -- (pitzer);
\draw [arrow] (pitzer) -- (calc_hr_gs);
\draw [arrow] (calc_hr_gs) -- (calc_hr_v);
\draw [arrow] (calc_hr_v) -- (calc_mixing);
\draw [arrow] (calc_mixing) -- (calc_excess);
\draw [arrow] (calc_excess) -| (converge);

% Flujo modelo simple
\draw [arrow] (simple_model) |- (converge);

% Flujo final
\draw [arrow] (converge) -- (output_sub);
\draw [arrow] (output_sub) -- (end_sub);

% === CAJA DEL SUBMÓDULO ===
\node[module_box, fit=(start_sub)(inputs_sub)(simple_model)(calc_excess)(end_sub), 
      label={[font=\Large\bfseries, purple!70]above:SUBMÓDULO: CÁLCULO DE $\Delta H_{\text{exceso}}$}] {};

\end{tikzpicture}
\caption{Diagrama del submódulo para cálculo de entalpía en exceso}
\label{fig:excess_enthalpy}
\end{figure}